\customchapter{RESUMEN}

El glioblastoma multiforme presenta una resistencia evolutiva que limita la duración del control logrado con temozolomida. Esta tesis propone GBMARL, un marco de aprendizaje por refuerzo multiagente adversarial que modela la terapia y la evolución tumoral como dos agentes que interactúan sobre un simulador de poblaciones sensibles y resistentes. La política de terapia se entrena con MAPPO bajo entrenamiento centralizado y ejecución descentralizada (CTDE), y se compara con IPPO, ausencia de tratamiento, dosis máxima tolerada y la heurística adaptativa de Gatenby.

El simulador se calibró integrando DepMap Public 26Q1 y GDSC2 release 8.5. El cruce recuperó 34 líneas de GBM con ensayos de sensibilidad y 24 con expresión completa. Para la temozolomida se obtuvo una brecha de sensibilidad aproximada de 9,1 veces entre las poblaciones sensible y resistente. La evaluación usa el tiempo a falla combinado, que termina cuando se supera el umbral de carga o cuando la fracción resistente alcanza la mayoría, registrando el modo de falla.

En la comparación nominal calibrada, quince políticas MAPPO alcanzaron una mediana de 34 días [27, 41] y superaron a Gatenby en 13/15 semillas; quince políticas IPPO alcanzaron una mediana de 31 días [23, 34] y superaron a Gatenby en 12/15. La diferencia pareada MAPPO--IPPO fue exploratoria ($p=0{,}013759$). El análisis de sensibilidad identificó una fragilidad importante frente a la reducción de $K$: con $K=0{,}8K_0$, la mediana de MAPPO cayó a 25 días. Por ello, los resultados constituyen evidencia \emph{in silico} de comportamiento adaptativo bajo el modelo especificado, no evidencia de eficacia clínica.

\noindent \textbf{Palabras clave:} glioblastoma; terapia adaptativa; resistencia evolutiva; aprendizaje por refuerzo multiagente; MAPPO-CTDE; temozolomida.
